\medskip

L'objectif de cet exercice est de déterminer la distance entre deux droites non coplanaires.

Par définition, la distance entre deux droites non coplanaires de l'espace, $(d_1)$ et $(d_2)$ est la longueur du segment [EF], où E et F sont des points appartenant respectivement à $(d_1)$ et à $(d_2)$ tels que la droite (EF) est orthogonale à $(d_1)$ et $(d_2)$.

L'espace est muni d'un repère orthonormé \Oijk.

Soit $(d_1)$ la droite passant par A$(1~;~2~;~-1)$ de vecteur directeur $\vect{u_1}\begin{pmatrix}1\\2\\0\end{pmatrix}$ et $(d_2)$ la droite dont une représentation paramétrique est : 
$\left\{\begin{array}{l c l}
x&=&0\\y &=& 1 + t\\z &=& 2 + t\\
\end{array}\right.,\: t  \in \R.$

\medskip

\begin{enumerate}
\item Donner une représentation paramétrique de la droite $(d_1)$
\item Démontrer que les droites $(d_1)$ et $(d_2)$ sont non coplanaires.
\item Soit $\mathcal{P}$ le plan passant par A et dirigé par les vecteurs non colinéaires $\vect{u_1}$ et $\vect{w}\begin{pmatrix}2\\-1\\1\\\end{pmatrix}$.

Justifier qu'une équation cartésienne du plan $\mathcal{P}$ est : $-2x +y +5z +5 = 0$.
\item
	\begin{enumerate}
		\item Sans chercher à calculer les coordonnées du point d'intersection, justifier que la droite
$(d_2)$ et le plan $\mathcal{P}$ sont sécants.
		\item On note F le point d'intersection de la droite $(d_2)$ et du plan $\mathcal{P}$.
		
Vérifier que le point F a pour coordonnées $\left(0~;~- \dfrac53~;~- \dfrac23\right)$.
\end{enumerate}
\end{enumerate}
		
Soit $(\delta)$ la droite passant par F et de vecteur directeur $\vect{w}$. On admet que les droites $(\delta)$ et $(d_1)$ sont sécantes en un point E de coordonnées $\left( -\dfrac23~;~- \dfrac43~;~-1\right)$.

	\begin{enumerate}[resume]
\item 
	\begin{enumerate}
		\item Justifier que la distance EF est la distance entre les droites $(d_1)$ et $(d_2)$. 
		\item Calculer la distance entre les droites $(d_1)$ et $(d_2)$.
	\end{enumerate}
\end{enumerate}
